

\newpage
\section{导论}

\newpage
\subsection{研究背景及意义}




\subsubsection{国际背景}
从国际背景看，当前全球经济正出现逆全球化趋势。截止2019年，美国和中国分别作为全球最大和第二大的贸易国。中美之间在贸易领域的任何波动，都可能影响世界经济的发展。中国近几十年的崛起，美国为首的西方阵营的一些国家对中国产生了许多恐惧，因此在贸易、金融、科技、外交军事等方面打压和威胁中国。2017年开始至今的中美贸易摩擦，不仅影响全球经济增长，也对全球金融业造成不可小觑的影响。2017年中国对美货物贸易顺差占中国货物贸易顺差的65.3%，位列第一，而以此同时美国对中国的货物贸易逆差也位列第一。为了扭转贸易逆差局面，美国开始用各种手段对华贸易摩擦。相关事件包括2017年以来对中国开展301调查、2018年美国对340亿美元的中国商品加征25%进口关税、2019年美方宣布美国企业不得使用对国家安全构成风险的企业所生产的电信设备，将华为公司列入出口管制“实体名单”。中美贸易摩擦使得金融市场，包括股票市场、债券市场、外汇市场均产生了不同的溢出效应。一方面通过影响市场基本面因素导致金融市场波动，另一方面通过影响投资情绪导致金融市场波动。
此外，公共危机事件加剧了金融危机的产生。2020年初的新型冠状病毒肺炎这类突发公共卫生事件在全球的蔓延，严重威胁世界各国的社会稳定，也对宏观经济产生了巨大冲击。随着我国国内疫情有效控制，复工复产，经济形势也逐步回向正规，然而国外疫情开始加重。疫情冲击通过降低需求市场和供给市场的偏好，沿着产业链、供应链、资金链等传导致金融市场，导致金融市场出现大幅震荡，全球各国股指自3月份起纷纷暴跌。包括美国股市在内的多国股市多次触发熔断。在经济低迷的情况下，利率一度低至新水平。许多国家采取量化宽松的极端货币政策，向市场注入流动性，刺激经济复苏。由于当前全球化较为紧密的环境下，一国的经济金融政策容易引发连锁反应，因此国际经济波动在一定程度上影响我国环境。

\subsubsection{国内背景}
我国经济经过几十年的高速增长，目前处于中高等收入国家，要迈入高收入国家的门槛，需要进行转型升级，重要的举措之一是拉动内需。近几年贸易摩擦和疫情等多重的外生冲击的复杂多变的国际背景下，2020年7月30日，中央政治局会议释放出“加快形成以国内大循环为主体、国内国际双循环相互促进的新发展格局”的信号。以国内大循环为主体，发挥中国超大规模市场优势，国内国际双循环相互促进，强调以国内经济循环为主不意味着关门封闭，而是通过发挥内需潜力，使国内市场和国际市场更好地联通、促进。这样一方面有助于遏制逆全球化发展，另一方面拉动内需，促进经济增长。习近平总书记在党的十九大提出要推动形成全面开放新格局，并多次强调包括实体经济和金融业在内的对外开放重大举措落地应宜早、宜快。
新冠疫情对我国经济与金融造成了冲击。我国银行业金融机构一直是金融体系的主体。从中国人民银行的统计数据可知，截至2019年末，我国银行业机构总资产占中国金融业机构总资产比例为91%，总负债占比92%，其规模远远大于证券业机构、保险业机构总资产之和。从银保监会的数据可知，2013年至今，我国银行业的总资产和总负债规模仍在大幅攀升，由此可见我国至今依然是银行主导的金融体系。我国银行业既受到疫情的明显影响，也是此次金融业抗击疫情、稳定社会经济的重要支撑。短期内银行业务将面临一定的负面冲击，使得信贷偏好风险下行。由于疫情管制等抗议措施的影响，企业、服务业等生产和消费减缓，经营压力增大，企业还款能力削弱，尤其对于中小微企业、三农企业信贷投放比例偏高的农村商业银行来说，其回收贷款的压力增大。为此我国相关政策正加大对受疫情影响较大的企业，尤其是对中小微企业的支持。
如果一国国内经济模式较为封闭，那么国际经济动荡对其影响甚微。但是我国是一个出口导向型国家，对外贸易依赖度约为1/3。特别是当前一个时期，中国正处于全面对外开放的新形势，中国开放的大门只会越开越大。因此，国际经济波动对我国环境的影响无疑是巨大的，不能忽视潜在的金融风险，尤其是系统性风险。习近平总书记在党的十九大报告中强调，要健全金融监管体系，守住不发生系统性金融风险的底线。在中共中央政治局第十三次集体学习时，习近平总书记更是强调防范化解金融风险，特别是防止发生系统性金融风险，是金融工作的根本性任务。



\subsubsection{研究意义}

本文中对我国银行业系统性风险的形成机制和防范与化解的研究具有较强的理论意义和现实意义。

理论意义：

关于我国银行业系统性风险的研究具有很强的理论意义。系统性风险的研究源于早期学者对传统金融危机的研究。历史上每一次的经济危机、金融危机都能促进经济理论、特别是金融危机理论的发展和完善。比如20世纪30年代的资本主义世界的大危机，引发了学界对于金融危机理论的研究热潮。20世纪80年代以后，随着金融全球化的深入和世界范围内的货币危机的频繁出现，关于金融危机的话题再次成为研究热点。在此期间，逐步形成了以Krugman为代表的第一代货币危机理论、Obstfeld为代表的第二代货币危机理论、以及众多研究者们依次发展和梳理逐步形成的第三代危机理论。1997年东南亚金融危机爆发以后，金融安全问题引起了广泛的关注。2008年金融危机掀起了全球对金融安全问题研究的新高潮，研究者们反思危机发生的原因，认为2008年金融危机主要是由次贷危机引发系统性风险造成的，由此巴塞尔银行监管委员会推出了基于宏观审慎监管角度的巴塞尔III协议，补充了宏观审慎监管角度对系统性风险防范的要求。对于当前背景下国内可能发生的银行业系统性风险的形成机制、包括防范和化解方面的研究正在一个新的势头银行业系统性风险还存在很大研究空间。
现有关于系统性风险的形成机制，包括传导路径在内，其研究包括关于金融恐慌与挤兑的流动性风险，银行机构关联性和传染机制等。本研究将对我国银行业系统性风险的形成机制和防范与化解措施紧密结合。从推动理论发展方面，以传统的金融恐慌与挤兑理论和银行体系的脆弱性理论和支付清算的理论为基础，探究当前一个时期，我国面对的国际冲击和国内压力下，实体经济危机如何引致信贷风险，进而引发系统性风险的形成机制，然后以经济金融周期波动理论为基础，设计在研究对象范围内相对合理的应对策略时，考虑我国银行业系统性风险如何反过来影响国民经济，作为评判其策略的优劣的标准之一。本文可以在一定程度上推进在当前我国新形势下的空白。从方法的改进和创新方面，现有研究关注于系统性风险的关联性和风险的放大和扩散机制方面，本研究基于此，深入细化银行间市场的网络分析方法，纳入动态随机一般均衡模型中，为扩展该模型做出一定的理论贡献。

实践意义：

银行业是目前我国金融业最重要的组成部分。银行业如果出现风险，必然影响整个金融业安全，乃至冲击实体经济。2018年中国银保监会推出的对外开放的新举措中，有一半与银行业直接相关。由此可以看出，中国政府对银行业对外开放越来越重视，开放程度越来越大。银行业近些年也出现了一些个别短暂的危机性事件，不应忽视。例如2019年5月的包商银行信用风险事件；2019年10月的河南伊川农商行挤兑事件；2019年12月的锦州银行重大资产重组事件等。虽然我国没有爆发过真正的银行危机，但是一旦爆发危机，对整个经济的冲击是深远的。为了从银行业层面防范与化解系统性风险，就应该分析银行业系统性风险的形成机制，分析主要成因，才能对症下药。与以往的金融危机相比，当前面对更为复杂多变的国际形势，和我国向高收入经济体迈进的转型需要、经济和金融体系越来越大地开放的双重趋势下，本研究通过分析银行风险的形成机制，以便于为依托制订防范化解系统性金融风险的最优政策，能够为我国银行业为主的金融业的稳定建言献策，实现不发生银行业系统性风险，实现“保六稳”中的稳金融，将具有十分重要的实践意义。


\newpage
\subsection{国内外研究综述}

综述部分的内容可以用三个维度组合归类：研究角度维、相关理论维、研究方法维。为了阐述方便，我们以研究角度维作为切入视角展开描述。该视角分为三个方面：银行业系统性金融风险（以下简称系统性风险）的形成机制（含系统性风险传导机制）、系统性风险测度、系统性风险的防范与化解。
在描述系统性风险之前，我们需要了解关于系统性风险的定义，才能更好地描述其形成机制。金融系统性风险是指一个或者几个重要金融机构的失败将通过金融机构之间的相互联系而引起其他金融机构的失败，并进而对实体经济经济产生实质性的负面效果。


\subsubsection{银行系统性风险测度相关的研究}

为了研究系统性风险形成机制，涉及如何测度系统性风险。在系统性风险的测度方法方面，主要涉划分为以下几类：基于风险组合模型法的测度、基于网络分析法的测度。

风险组合模型法也称为简约化方法，该方法由衡量证券组合风险的方法发展而来，将整个系统看成是金融机构的组合来衡量系统性风险。其特点是，主要基于金融市场数据估计金融机构之间的某种相关性，然后根据相关性来评估风险的传染性。该方法的优点是可以识别组成系统的各机构的共同风险因素，跟踪某一机构出现问题将如何影响其他机构，度量单个金融机构对整个系统的风险贡献。该类方法主要的方法有VaR、CoVaR、GARCH、期望边际损失MES、或有权益分析方法CCA等。(Engle, 2000)提出的DCC-GARCH模型应用广泛。(白雪梅等, 2014)采用CoVaR方法分析了我国金融系统的稳定性，并建立了一个预测系统性风险的模型，该研究认为CoVaR方法相比VaR方法，可以更好测度尾部损失，而现实中的金融数据往往并非正态分布，而是呈现尖峰厚尾分布。CoVaR的缺点在于其采用相关性分析反映风险的外溢效应，难以分析因果性。且其危机时期和非危机时期的表现存在很大差异。风险组合模型法的缺点是忽视了风险传染渠道，对风险传染过程难以刻画，在反映机构关联性方面存在明显不足。

为了更多测度与分析金融系统之间的关联性，网络分析法通过研究金融系统之间的关联性这一重要因素研究系统性风险的成因、系统性风险的传导途径、系统的脆弱性等。

网络分析法分为两种研究视角，其中第一类走仿真路线，即从微观层面研究银行间市场的构成的网络的流动性、稳定性、抗毁性。这类研究很多结合第一类视角的复杂网络结构，基于经济学理论，构建相关的模型，得到算法，利用模拟仿真的方法做压力测试实验。在系统性风险的形成机制中，我们已经罗列过相关研究。第二类走统计实证路线，也就是从宏观层面测度系统性风险的路线。研究银行间市场构成的复杂网络的结构的特征，比如用金融机构之间的相互关联性和集中度等复杂网络统计领域的指标测度。国外的学者通过获取德国、意大利、瑞士、奥地利、墨西哥等银行的同业拆借数据对银行间市场进行测度，研究银行关联性对于银行体系稳定与风险传染的关系。例如(Boss等, 2013)分析了奥地利银行同业拆借市场，通过观察3年间900家银行的数据，结果显示银行间同业拆借市场内在的联系分布服从幂率指数为1.85的分布。(Martinez-Jaramillo等, 2014)利用墨西哥银行双边同业数据进行实证研究，发现该国银行具有无标度特征。(彭建刚等, 2013)基于(Allen等, 2000)的研究的四种代表性的银行间市场网络的结构的机理，分析了我国银行同业交易商，分为四大类，分析了拆借市场的交易主体的结构的变化。由于中国银行间同业拆借中心并没有公布分类交易商之间的双边数据，根据分类交易商的拆入拆出数据，采用最大熵方法大致估计参数作为同业市场结果的基本特征。发现我国同业市场结构已经由单向的资金中心结构向类似于完全交易市场结构转变。(Gai等, 2011)研究银行同业拆借网络更大的复杂性和集中度如何放大了系统的脆弱性。研究发现：更严格的流动性监管和宏观审慎政策可能会提高系统的弹性；网络的复杂性和集中度对于确定财务传染至关重要；银行间市场细分的网络模型突显了回购市场活动的作用。在系统性风险的指标创新方面，(Ouyang等, 2020)考虑了投资者注意力对金融机构风险传染的影响以及金融网络结构所带来的潜在风险，因此使用金融机构的百度搜索索引来描述投资者的情绪，采用linear quantile lasso regression 和 local polynomial method 方法估计TENET模型，用有向网络方法，衡量中国银行业的系统性风险蔓延效应。研究发现：当前银行具有的“太大而不能倒闭”和“过于联系而倒闭”的特征；在整个样本期内，整个网络的连通性和系统性风险具有相同的趋势，百度搜索指数会放大系统性金融风险。在网络分析法中对于系统性风险的测度指标中(方意, 2016a)采用SR、VBI、CBI等指标，由于银行体系总溢出损失。受银行体系初始资本金规模大小影响。SR指标剔除了银行初始资本金规模对系统性风险大小的影响。更能准确地反映银行体系的系统性风险水平。(方意, 2017)，在研究中通过纳入去杠杆——降价抛售机制，进而对中国最为现实的几类冲击（房地产贷款违约、地方政府融资平台贷款违约）进行压力测试，并利用银行倒闭冲击量化地论证了本网络模型相对于传统网络模型的优势。

在系统性金融风险中，很多时候，如果能够识别一些具有风险特征，可能对整个系统造成很大风险影响的金融机构，则能够有利于防范与化解系统性风险。这样的机构成称为关键性的金融机构（SIFIs）。在识别关键性金融机构的过程中，可以采用压力测试方法。压力测试是一系列用来评估一些极端但又可信的宏观经济冲击对金融体系脆弱性影响的技术总称，而宏观压力测试由于能模拟潜在金融危机等极端事件对金融系统稳定性的影响，引起了国际监管组织和各国监管当局的格外重视。各国的实践中压力测试往往包含几个步骤：一、确定需要测试的金融机构或资产组合的范围；二、压力情景的设定及强度选择；三、量化分析宏观压力请将对银行系统的冲击；四、根据测试结果评估银行体系风险承受能力。各国央行普遍采用压力测试法进行研究。我国的实践中，中国人民银行近几年来通过压力测试方法中国银行体系进行压力测试。在学术界中，压力测试法也得到广泛使用。在系统性风险的形成机制中，我们已经提及过相关研究。基于网络的压力测试方法将网络分析方法纳入到压力测试方法的框架中。相关研究有很多，例如(方意, 2017)通过纳入去杠杆——降价抛售机制，进而对中国最为现实的几类冲击（房地产贷款违约、地方政府融资平台贷款违约）进行压力测试，并利用银行倒闭冲击量化地论证了本网络模型相对于传统网络模型的优势。


\subsubsection{银行系统性风险的形成机制相关的研究}

关于系统性风险的形成机制，应当从系统性风险的理论来考察。银行系统性风险定义为由于金融体系内机构倒闭、信用风险、金融恐慌造成的挤兑现象、资产波动、实体经济冲击等一系列极端外部冲击造成银行的流动性风险，并且经由一个机构传染到互相关联的多个机构，最终对实体经济产生负外部溢出的风险。尽管有很多学者提出了各种导致系统性风险的成因，但是归纳现有梳理的文献，可以按照不同的角度分类成因。总的来说，可以分为以下几个要素：冲击、一致性、杠杆性、关联性。其中，冲击包括外生冲击和内生冲击。外生冲击包括诸如经济实体部门的违约现象对金融部门造成的冲击、金融机构持有资产价格变动、国际金融市场冲击、经济周期因素、公共突发事件造成的影响、政策带来的影响等。内生冲击主要包括金融的脆弱性、机构关联性、机构行为的同质化、道德风险等。此外，政策的反馈效应也被认为是内生冲击的一部分。所有的内生冲击最终都可以溯源到外生冲击。因此，冲击在此特指外生冲击。一致性主要指储户、投资者、企业等部门获取信息、结构组成、行为方面的一致性，以及银行机构部门个体表现出来的在规模、资产负债结构、决策风格等方面的一致性。杠杆性主要指银行部门存在的高负债特性。关联性主要包括几方面：一是银行在金融市场资产的关联性，由此表现出相似性；二是银行间市场资产负债的关联性；三是银行业务之间的关联性。银行间市场的关联性如同硬币的两面，一方面可以为银行间的资金流动提供畅通的渠道；另一方面，也使银行系统面临更高的传染风险。

由上述几个系统性风险的形成机制要素的相互作用可以将系统性风险分为以下研究：

在外部冲击方面，如国内外背景所述，当前我国系统性风险的重要外部成因之一是宏观经济波动造成的冲击。当中很重要的冲击包括突发事件造成的冲击。例如(和文佳等, 2019)量化分析了中美贸易摩擦对中国银行业系统性风险影响的水平效应和趋势效应。(杨子晖等, 2020) 如量化分析了“新冠”疫情等突发公共事件对我国宏观经济与金融市场的冲击影响。

在冲击与一致性共同作用方面，有储户与投资者挤兑冲击与群体行为一致性因素的作用，相关研究主要是银行间预期协调研究。该研究主要起源于(Dybvig, 1983)从挤兑的角度对提出了D-D挤兑模型作为该研究的开创。D-D模型，主要研究银行挤兑如何发生，如何引起银行破产。提出的挤兑发生的思想是，设定两类家庭，短期消费家庭取出存款进行消费的时间会早一点，长期消费家庭有的家庭会迟一点。银行会将这些存款投资于2期，且到第2期才会获得投资收益。短期消费家庭取回的是他们最初的存款，不含收益。长期消费家庭会获得正的投资收益。银行面临的问题是，当他们最初进行投资时，他们不知道有多少储户属于短期消费家庭，因此他们不确定自己是否有足够的流动性资产来应对短期消费家庭的取款需求。一旦银行低估了这个数量，那么除非银行可以从银行间市场借到资金，否则会缺少现金。此时当储户希望自己能全额取出自己的存款时挤兑就发生了。研究提出的观点是，为了避免这种情况的发生，需要由政府提供存款保险以规避危机的发生。后续研究者主要针对存款人与银行间挤兑时间发生的概率的预期，研究银行负面信息策略性披露与银行挤兑的关系，由于这个研究角度难以在实践中获得数据，因此大多采用理论建模的角度或者实验仿真的方法进行，相关研究例如(Garratt等, 2009)。

在银行一致性和关联性结合的研究角度看来，(Acharya and Yorulmazer,2008)、(Benoit et al,2017)提出观点认为金融机构之间，一方面通过相互融资和业务交易，以共同抵御流动性冲击，实现风险分担。另一方面这本质上体现为金融机构的流动性风险管理。在冲击对一致性的影响的研究角度来看，有资产相似性的研究。其中包括，由于银行体系持有相似的资产，当金融市场的资产价格下降时，或者当外部实体经济整体形势恶化时，银行体系面临的整体冲击。相关研究有(Acharya, 2009)基于DSGE模型，研究各银行资产配置区域同质化，在冲击时的一致性。而从冲击对杠杆性的影响的研究角度来看，一般认为，当银行受到资产价格变动或者实体经济造成的初始冲击时，银行体系由于高负债的杠杆性，导致初始冲击被放大。而当关联性引入到其中，则会看到，放大的初始冲击在具有关联性的银行间市场传染，二者互相正反馈造成恶性循环的局面。

冲击对银行间市场的关联性存在影响。首先，银行间市场显现出明显的关联性要素。在银行间市场中，银行间通过借贷、支付清算、贴现、承兑、担保等形式形成了错综复杂的关系，呈现出高度的复杂度、集中度、关联度。风险在银行间市场之间的传染机制有一些研究。相关的研究方法主要有网络分析法、压力测试、仿真模拟方法，或者数理分析中的随机过程、博弈论等方法。(Allen等, 2000)的早期研究，从银行间市场的关联性的角度，创新性地将风险在银行间市场的分担问题进行了描述，奠定了该领域微观经济学基础。其提出的模型构建了三种包含四个银行的网络结构单元。发现如果银行间市场结构是不完全市场，虽然可以实现最优风险分担，但在面临流动性冲击时更容易出现风险传染，从而使得系统更脆弱；如果银行间市场结构是完全市场，那么可以实现最优风险分担，可以弱化某个银行的初始流动冲击给系统带来的影响。后续相关的研究包括冲击对银行间市场的影响，包括冲击规模、冲击目标、银行规模、网络结构、银行关联度等对于银行破产比例的影响。例如，(Lorenz等, 2009)分析银行破产数量与规模及其概率与银行间关联度之间的关系，研究发现二者存在“U”型关系，即当银行间中等关联程度时，银行破产数量与规模会最小。(李守伟等, 2011)从复杂网络视角，通过删除网络的节点和边以分析网络的连通性，以此方法模拟银行倒闭，研究银行网络结构在面对随机性攻击和选择性攻击的承受能力来分析银行间市场的稳定性，研究发现银行网络对于随机性攻击具有较高的稳定性，对于选择性攻击具有较低的稳定性，并且两种攻攻击方式对银行网络影响的差异大小与阈值有关。(李守伟等, 2012)基于(Allen等, 2000)，分析随机性冲击和选择性冲击下传染的风险特征和稳定性。通过设定银行间市场网络结构为随机网络、小世界网络、无标度网络，给出构建方法。得到模拟分析结果：冲击在小世界网络中造成的传染风险效应最大，在无标度网络中最小。意味着在此三种网络结构下，银行间市场无标度网络面对冲击具有最高的稳定性。(张锋华等, 2013)在银行倒闭诸多原因中，流动性危机比资不抵债对银行的影响更大。为模拟金融危机中银行网络流动性再分配过程，以研究银行网络的流动性转移效率，结合复杂网络理论与(Eboli, 2010)的流动性转移模型，在有向有权随机网络的基础上构建了一个银行网络流动性动态配置的多主体模型。通过向银行网络施加流动性冲击来模拟金融危机中银行网络流动性再分配过程。模拟结果表明银行网络结构、连接度对网络流动性转移效率有重要影响。同等条件下，银行网络流动性转移效率与流动性冲击的大小呈负相关，与网络结构均匀程度及连接度正相关。为此，作者提出两种有效手段用以提高银行网络流动性转移效率：一是增强银行间的信贷联系，并力求同时让银行网络更加均匀；二是减少银行网络受到的流动性冲击；相应的政策建议对应于完善我国银行间拆借市场、充分发挥央行的最后贷款人角色。(彭建刚等, 2013)基于(Allen等, 2000)的研究的四种代表性的银行间市场网络的结构的机理，采用压力测试的方法对我国银行间同业市场上分类交易商之间的流动性和风险传染效应进行了研究。研究发现小规模流动性冲击下，风险具有收敛性，但是在大规模流动性冲击下会导致流动性风险蔓延和同业市场交易量萎缩。(Gai等, 2010)探讨了网络结构的变化和资产市场流动性如何影响传染的可能性和潜在影响。通过构建具有任意结构的金融网络中传染的分析模型，发现金融系统表现出了强大的但仍脆弱的趋势。尽管传染的可能性可能很小，但是当问题发生时，这种影响可能会非常广泛。(Ladley, 2013)比较了大规模冲击和小规模冲击下银行间市场的作用。研究发现，一般而言，在大规模风险冲击下，银行间市场起到了放大器作用；而在小规模冲击下，银行间市场则主要起到风险分担和稳定器的作用。(Marcel等, 2014)研究认为，金融机构关联是系统性风险的一个主要驱动力，在任何系统性风险分析中，银行间同业拆借都值得高度关注。

此外，在银行间同业网络的偿付清算问题的研究中，(Eisenberg等, 2001)开创性地描述了复杂金融清算系统的特性，其应该满足的条件，定义了清算向量，为其存在性和唯一性提供条件。其证明了存在一个唯一的清算支付向量，用以满足清算系统的义务。通过分析清算向量的性质，提供描述清算向量与金融系统基础参数之间关系的比较静态，较早提出和解决了简单的同业网络中的支付清算问题。(Fouque等, 2013)提出了一个简单的银行间借贷模型，其中各银行的对数货币储备的演化用漂移耦合的扩散过程系统来描述。为了更多地将模型扩展和细化，研究者们以此为基础做了研究。(Carmona等, 2015)提出了一个带有博弈特征的银行间借贷模型。各银行的货币储备的演化用漂移耦合的扩散过程来描述，各银行通过控制向中央银行的借贷利率来实现最优化。研究的主要方法是基于平均场博弈，构建精确的纳什均衡解。先用Pontrayagin随机最大定律获得FBSDEs的解，以构建模型的开环均衡。再计算闭环的Markov均衡，这其中分别通过两种方法实现：一是用改进的在开环的情形中的Pontryagin随机最大化定律实现；二是基于用动态规划方法解HJB偏微分方程来实现；此外讨论了流动性项的金融含义和央行的角色。研究的结论是，银行间借贷为系统创造了稳定性，模型中，央行表现为票据清算中心的角色，其提供额外流动性的行为增强了这种稳定性，以抵御系统性风险。进一步地，(Carmona等, 2016)通过引入延迟，讨论了没有延迟因素和有延迟因素的差异。研究发现，由于控制存在延迟，需要归还或收取退款而产生的新效果降低了在没有延迟的情况下观察到的流动性。(Banerjee等, 2018)结合了基于静态分析的复杂网络方法，和基于动态分析的平均场理论，扩展了(Eisenberg等, 2001)的离散时间的状态至连续时间的状态，以及将单期的偿付过程扩展至多期。
在由冲击、杠杆性与关联性共同作用的研究方面，主要在于初始冲击使得银行折价抛售资产获取流动性，在杠杆性的作用下放大了冲击，使得银行陷入更危险境地，为此银行赎回其它与该银行相关联的其它银行的资产，造成次生冲击并在具有关联性的银行间市场传染，二者互相正反馈造成恶性循环。相关的研究有(Diamond等, 2005)。(Cifuentes等, 2005)通过仿真模拟的方法研究了资本充足要求导致的资产价格变化的流动性风险传染渠道。在一个相互关联的金融机构市场中，针对流动性风险传染的两种渠道（直接互联的金融机构资产负债表、和资产价格变化渠道）进行了对比研究。这两种渠道直接通过银行间市场与资产市场这两个市场，基于正反馈机制形式将这种流动性风险不断传播并加大，前者被称为直接传染渠道，后者则被称为间接传染渠道。在国内，(方意, 2016a) 认为我国的银行业的风险在于房地产业等借贷规模巨大的风险，加上经济下行压力，容易对银行业造成冲击，而密切关联的银行业容易导致风险的扩散和累积。在模型上创新性地构建了包含银行破产机制和去杠杆机制的资产负债表直接关联网络模型，通过四类银行特征和四类外生参数影响的四类传染渠道，度量了系统性风险，并且对其传染渠道进行了分析，通过引入最低杠杆率监管要求，分析当遭受损失冲击时，银行业间的风险传播作用，以及为了满足最低杠杆率监管要求从而抛售资产带来的损失，导致杠杆更大，从而加剧风险。
此外，在由冲击、一致性与关联性共同作用的研究方面，(Manz, 2010)通过一个基于信息的危机传染全局博弈（Global Game）均衡模型提出了“正向”传染的概念，即一个金融机构在财务上的失败可以诱发另一家金融机构的失败，且如果这两家机构基本面信息相关性越强，则传染效应越强。


\subsubsection{银行系统性风险的防范与救助的研究}

银行业系统性风险防范与化解方面的研究，涉及到多种政策工具或者政策工具组合的运用。常用的政策救助工具包括：杠杆类（资本类）工具、流动性工具、关联性政策工具等。杠杆类（资本类）工具，其特点为关注银行的资产负债结构，工具包括资本充足率、资产/权益金比例。杠杆类工具的优点是，可操作性强，政策成本低，缺点是存在顺周期性。巴塞尔协议III的宏观审慎逆周期的提出就是为了改进杠杆类工具的不足。常见的杠杆类政策工具是注资政策，其操作行为是监管当局直接向银行机构注入资金，使得银行机构能够在一定时期内获得资本金。流动性工具则关注银行与资产负债的流动性问题，包括流动性覆盖率、净稳定资金比率、流动性比例等，其优点是可以动态监测银行的流动性状况，主要操作是政策救助部门利用救助资金规模购买银行的非流动性资产，使得受到救助的金融机构获得流动性资产。在这其中涉及到商业银行对流动性资产的抛售顺序的决策，也就是关联性的政策工具。关联性政策工具通过改变银行机构之间的关联性，从而调控系统性风险。例如监管当局用银行拆分和合并、窗口指导等方式。政策救助部门利用救助资金购买银行抛售的非流动性资产，可以改变银行机构之间的关联性。(方意等, 2019)在资产抛售与监管方面，分析不同抛售规则下银行体系系统性风险的差异，借助博弈模型对比分析微观和宏观审慎政策下银行机构的行为选择差异性，刻画宏观审慎政策对系统性风险的传导机制。从而测算政策实施的收益与成本，研究发现资产抛售规则与监管方面，先抛售流动性资产相对于先抛售非流动性资产可显著降低银行体系系统性风险。然而，在实践过程中，若不存在监管部门强制规定。各银行不会选择先抛售流动性资产。通过宏观审慎政策中的窗口指导政策，或者流动性注入政策，当局通过奖励和惩罚来调控银行的资产抛售顺序的意愿，从而改变银行自身的资产抛售规则。

由于早期DSGE框架侧重于凯恩斯经济学的成分，较多分析消费、生产等领域，而对银行部门的考虑较少。很多家庭或者企业储蓄的细节作为外生变量，其具体细节较少被考虑。中央银行发挥救助银行体系的职能，向银行体系注入流动性的机理应该被给予更多的阐述，因此DSGE框架下研究系统性风险和金融危机，成为一个研究趋势。DSGE框架越来越多被世界各国央行作为宏观经济预测和政策分析的工具，比如(Rostagno等, 2010)开发的欧洲中央银行相关模型；IMF的(Benes等, 2014a, 2014b)基于DSGE思想，构建了MAPMOD的理论结构，研究与过度信贷扩张相关的漏洞，并进行宏观审慎政策分析。

在诸如挤兑的违约风险，以及在银行间市场的关联性方面的传染的研究方面，这些年的研究中，国内外已有学者开始将银行部门、金融加速器以及银行间市场和银行资本等多重金融因素，融合到DSGE模型中。(Allen等, 2009)建立了一个分析银行间市场以及央行应如何干预的简单的两周期理论框架。银行可以持有具有较高收益率的一期流动资产或二期长期资产。研究结果表明通过中央银行的从事公开市场操作可以解决银行间市场的纯粹的流动性风险。(Miller等, 2010)也研究了类似的问题。为了维持银行体系流动性，降低系统性风险，传统的货币政策在利率的零下限约束状态下可能难以起到更好地作用，因此需要非传统的货币政策加持。在DSGE模型在金融中介的研究方面，(Gertler等, 2012)等人在金融危机之后研究了包含金融中介的DSGE模型，该模型几种研究流动性风险、政府干预性的政策、以及对资产回报风险的认知等是如何影响金融中介的。该模型的金融摩擦和扰动源是银行可以因为私人的或者非生产性的动机而挪用它们从银行间市场获得的资金。但是该研究没有涉及到关于违约的这一类重要问题。(Cúrdia等, 2011)在运用DSGE模型通过讨论中央银行的资产负债表阐述了非传统货币政策的作用。里面引入了违约相关的变量，但是是假设作为外生的作用。为了将违约相关的变量内生化，(Wickens, 2012)阐述了一个违约风险是内生定价的银行模型，将DSGE框架纳入其中。(Gertler等, 2015)融合了金融加速器效应和D-D银行挤兑模型，构建了DSGE模型，给出了金融机构部门与宏观经济之间在下行周期的共振机制。当宏观经济疲软时，产出和资产价格的下降会降低金融机构的盈利能力，从而其资本金遭受损失。此时，储蓄者预期金融机构偿付存款能力下降、信息不对称程度增强，会对金融机构进行挤兑，从而使得金融机构的资产负债表快速收缩。信贷资产的萎缩会提高商业信贷的价格，使得实体经济进一步萧条，加深金融机构资本金和资产价格的下降幅度。随着金融机构资产负债表状况和资产清算价格的不断恶化，金融机构的破产概率会随之上升，最终形成“宏观——金融”间的恶性循环。在把银行间市场纳入到DSGE框架的方法中，(王擎等, 2016) 基于中国经济金融特点，构建了DSGE模型，分析不同的因素对系统性风险的影响，模型包含居民、银行、企业、央行（含政府部门）四大部门。其中银行部门和居民部门按照资金富余程度进行了细分，通过不同的银行部门之间的拆借行为引入银行间市场。该研究使监管政策的研究跳出难以定论的微观风险效应视角，从宏观审慎层面分析和研判资本监管政策的有效性和改进方向。但是该研究的缺点在于对于居民和银行间市场的划分只对整体采用固定的比例划分，没有考虑居民之间、银行之间可以在不同类型转换的细节。由于我国银行业的特殊性，未考虑到当挤兑状态时，无法偿还居民存款的情况。

现有的一些针对国内的研究主要局限于单个或几个因素、或者几个政策工具的作用。(方意, 2016b)在外生杠杆模型的GNSS模型(Gerali等, 2010)的基础上，引入了中国特色的金融监管工具——存贷比监管工具，通过构建DSGE模型，同时研究了三种宏观审慎政策（贷款价值比、资本充足率、存贷比）及其组合，对于三类金融稳定目标（实际产出、信贷、房地产价格）的有效性。在一个模型框架中研究政策工具的组合对于多目标的影响，不仅更体现出实际经济运行的复杂性，也符合央行的研究方向。在银行业系统性风险防范与化解方面的研究中，当央行和各商业银行部门同时追求多个目标时，其面对的就是一个多目标的优化问题。这本身在经济学中是很常见的问题。各部门可以理解为在一个动态随机一般均衡模型（DSGE）中，构造各个目标的线性函数或者对数线性函数形式的目标函数，通过选择合理的政策变量，来最优化这个总的目标函数(周小川, 2016)。在多目标优化过程中，银行业风险对经济的负面影响应该作为监管当局和政府部门在采取决策时衡量的目标之一。比如说，在银行业冲击对经济周期的影响(Dib, 2010)研究了银行业在总冲击的真实影响的传递和传播中的作用，评估了金融冲击在美国经济周期波动中的重要性。研究通过构建将活跃的银行业纳入到具有金融加速器的DSGE模型中。

\subsubsection{文献评述}

以上罗列系统性风险的形成机制、测度、防范与化解相关的研究。这三个方面之间的关系在于，系统性风险的测度侧重从宏观角度构建相关指标，衡量系统性风险程度。形成机制侧重从微观角度分析系统性风险的成因、传导机制。二者互相提供实证支持和机理支持。系统性风险的防范与化解则基于形成机制与测度分析，寻求防范与化解系统性风险的对策。现有的测度分析能够通过数据和相关测度指标、方法，识别重要性金融机构。

这些形成机制的模型的优点是考虑了实际金融机构的关联性而能对系统性风险在微观层面上做较好的刻画，但是缺点也有，一方面从测度的角度来说，往往考虑截面数据，而对于时间维度考虑不足，另一方面从建模成分和国情来说，大多数研究考虑的传染因素过于单一，仅包含银行破产这一传染因素，与现实运作还有很大差距，银行破产在我国尚不完善的银行破产机制中并不多发生。截止2019年末，我国仅有海南发展银行等少数几家银行出现破产现象。因此，至少在当前时期，按照国情来说，我们考虑的银行破产因素实质上应该是技术性破产，即陷入困难，而不是实际意义上的破产，从而进入破产重组程序。

中国人民银行近几年来通过压力测试方法中国银行体系进行压力测试。但是，这些对中国银行业进行压力测试的权威报告只考虑资产面冲击导致的直接损失，对银行网络体系相互作用导致传染性损失研究甚少。但是该方法由于其没有考虑银行之间的关联性造成的传染风险，低估了不利宏观经济冲击导致的银行总风险度的系统性风险。因此建议结合压力测试法与网络分析法进行分析。

对于系统性风险的银行体系的关联性相关的研究中，无论是在复杂网络方法方面进行实证研究，或者是构建模型的算法进行模拟仿真研究，虽然在一定程度上接近对金融系统的真实结构的描述，但是有一个比较大的问题是，其较多局限于金融部门内体系、或者具体的金融部门间，在一定程度上忽略了经济系统领域的其它部门参与作为一个实际的经济整体而对金融部门的相互作用。所以我们需要一个即能够兼顾金融部门体系与宏观经济体系，又要兼顾理论与现实的框架。但是在这个过程中，也应该考虑研究的重点，在整体框架的把握上，做到在主要研究对象上进行细致深入的刻画，在于主要研究对象业务相关联的部门的刻画上做到恰到好处。央行在对系统性风险防范和化解的政策中，更多会将其纳入到宏观经济体系中考虑。因此，DSGE分析框架适合作为研究系统性风险的工具。此外，我们的研究中将考虑对多目标优化的综合衡量。

但是在各国央行使用DSGE分析框架研究系统性风险与宏观政策相关的研究中，可能存在的缺点在于未能细分金融机构，尤其在引入银行间市场层面较为简单，未能更详细刻画各金融部门的相互作用。因此，我们考虑将复杂网络分析方法纳入到DSGE分析框架中。

综合上述研究，本文拟在上述研究的基础上，以我国的银行业的系统性风险的成因为研究对象，以融合银行间市场网络模型的DSGE框架为模型的基础，更全面地研究：第一步，初步筛选几个重要的可能导致系统性风险的成因的影响因素，然后使用模型框架，分析我国银行业为主的系统性金融风险的主要成因，以及第二步，以其形成机制为依托，分析政策工具的适用性，研究在系统性风险发生时，采用什么样的政策工具组合，可以实现风险最小化，同时兼顾经济的平稳运行。通过现有研究，发现系统性风险当中的刻画银行间市场的方法，在DSGE框架中往往过于简化，以至于难以刻画其实质。因此研究中我们将细致考虑如何在宏观计量模型框架中引入银行间市场网络。此外，对于我国银行业的现实国情，由于商业银行实质性破产的情况极少出现，因此在目前的研究中，我们暂时不考虑国外研究中，银行已出现资不抵债就破产退出分析的情况，而是考虑技术性破产作为的替代。技术性破产中，政府往往会通过一系列措施对银行施以援手。此外，我们在研究中将不对各类政策做诸如宏观政策、微观政策等概念上的严格分类，而是更注重政策的细节及其组合分析。


\newpage
\subsection{研究内容和方法}

\subsubsection{研究内容}
本文将通过构建以我国银行业为主的系统性金融风险的框架，研究系统性风险成因、不同的政策工具救助对系统性风险的影响，从而对防范与化解我国的银行业系统性金融风险提供政策建议。

主要研究如下：

研究的思路

主线思路：

本研究以银行关联性和挤兑理论为基础，构建系统性风险的度量的指标，结合相应的模型，探讨银行间市场在宏观经济体系中，面临冲击时，对系统性风险的影响，从而判定相应的系统性风险主要成因。以及研究在宏观层面，央行等监管当局针对系统性风险的主要成因，为满足整个金融体系稳定的目标，应当如何采取最优的政策工具的组合作为救助策略。

研究中拟采用的主要分析框架：

根据现有研究，我们的研究也将考虑银行挤兑导致的风险传染这一角度，探讨银行与银行之间如何通过措施加强银行间市场稳定性。我们拟重点研究由于违约等几类冲击，造成银行业风险传染的影响因素对于系统性风险的主要成因分析。分析过程拟采用基于网络模型的压力测试方法，以整合网络分析法和压力测试法的优势。

在对金融系统性风险的防范与化解的研究方法中，近几十年来，结合宏微观、计量经济理论于一体的动态随机一般均衡模型（Dynamic Stochastic General Equilibrium Model,简称DSGE）逐渐成为经济分析工具的明星。DSGE并非某个具体的模型，而是一类分析经济理论的工具的思想体系和方法。在现有的研究中，DSGE模型框架大多是从新凯恩斯模型的总需求方程、价格调整方程、货币政策规则三个方程为雏形，通过引入各种部门、扭曲、摩擦等因素进化而来。DSGE模型在宏观经济研究和政策分析领域发展速度快、技术化程度高、应用范围广。其模型吸收了理性预期、动态优化、一般均衡分析的这些近现代经济学的革命性发展成果，可以分析经济增长、经济波动、货币与财政政策。与传统的计量经济模型相比，传统的计量经济模型在模型设定时并不是严格依据理论上得到的行为方程，只是利用了最终得到的变量之间的相互关系，从而在模型设定上具有一定的任意性，这样容易产生Lucas批判问题，这将对政策分析和评价造成很大影响，并且传统的计量经济模型并没有对经济的稳态进行显性的刻画，而DSGE模型可以避规此类问题。DSGE的特点在于动态、随机、均衡三个因素。首先，宏观经济现象是动态的。多种行为主体相互联系，在产品市场、劳动力市场、金融市场中做出跨期决策行为，因此静态模型都无法很好地描述复杂的经济现象。其次,宏观经济充满了各种不确定性。经济本身或行为主体本身都面临着各种内部或外部不确定性冲击,随机模型能很好地对这些不确定性进行建模。最后,宏观经济现象是所有经济行为主体相互联系、共同作用的结果。局部均衡无法描述所有主体动态作用的结果, 因此需要在一般均衡框架下进行分析。DSGE分析框架在研究与实施政策组合方面具有天然的优势，这些组合包括货币政策、财政政策、审慎政策等。



子研究一：

研究目标：构建我国银行业系统性金融风险的测度

研究思路：通过我国国情、现有文献进行分析，提出现有的可能的银行业系统性风险的主要影响因素、选取和构建测度系统性风险影响因素的中介指标，构建测度系统性风险的最终目标，用于测度系统性风险指标。指标部分，拟参考中国银保监会和现有一些文献中常用的监管工具实施的指标。系统性风险最终指标的用途主要是两个：一是在后续的形成机制的分析中，用于测度系统性风险在不同成因中的大小；二是在后续的分析防范与化解系统性风险中，用于测度系统性风险在面对不同的救助措施时期系统性风险的大小，查看是否有显著下降；

子研究二：

研究目标：分析我国银行业系统性金融风险的主要形成机制。

研究思路：通过我国国情、现有文献进行分析，提出现有的可能的银行业系统性风险的主要影响因素、选取和构建测度系统性风险影响因素的中介指标，构建测度系统性风险的最终目标，用于测度系统性风险指标，构筑宏观计量经济模型研究框架进行分析。指标部分，拟参考中国银保监会和现有一些文献中常用的监管工具实施的指标。

在对国情的把握中，例如，在分析当前我国银行业系统性风险的主要形成原因方面，我们参考2019年人民银行的分析报告。2019年上半年，人民银行选取了1171家银行开展压力测试，对其中资产规模在8000亿元以上的30家大中型商业银行开展偿付能力、流动性风险压力测试，对其余1141家银行开展偿付能力敏感性压力测试和流动性风险压力测试。评估银行体系在“极端但可能”冲击下的稳健性状况。测试内容包括：偿付能力压力测试，考察宏观经济下行、整体及重点领域风险状况恶化对银行资本充足水平的不利影响；流动性风险压力测试，考察政策因素、宏观经济因素、突发因素等多种流动性风险压力因素对银行各到期期限的现金流缺口的影响。测试结果表明，参与测试银行整体抗冲击能力较强；信用风险是测试银行的主要风险来源；市场风险对测试银行的影响有限；央行认为，银行体系流动性风险承压能力需进一步增强。由此可见，我们可以考虑纳入信用风险到当前系统性风险的成因之一。

参数估计部分，拟考虑利用部分公开数据等作为参考（例如我们将寻找国内外一些权威文献的数据、交易量数据可以考虑BankScope、国泰君安、wind等数据库、利率数据考虑Shibor等、各上市银行的详细数据从银行的年报中获取等等），对参数进行校准和估计。估计方法拟采用较为有效的Bayes参数估计方法。

模型构建和定量分析方面，对宏观经济系统的各部门划分与功能构建。我们对研究的重点对象，即商业银行部门，更多考虑银行与银行之间的宏观关联性与银行的微观异质性的特点。例如，按照规模分为大型银行和小型银行；每个商业银行之间都有各自最优目标；商业银行之间存在竞争和合作的行为等。

在分析违约造成的因素导致的流动性短缺时，继而造成危机传染的相关影响因素时，考虑构建银行间市场的风险传染网络模型。通过基于银行间网络压力测试法，判断银行间市场的关联性等因素对系统性风险的影响。我们拟将基于银行间网络子模型纳入到银行业风险传染子模型的银行系统性风险的宏观计量经济主模型框架，即DSGE模型中，通过经济系统中的冲击因素的变化，对银行体系做冲击响应分析，从而得到相应的结果，进而计算系统性风险的最终指标，判断各影响因素对于系统性风险的影响大小，从而判断我国银行系统性风险的主要成因。此外，我们可能会考虑使用SVAR等计量方法对上述的模型的判断结果进行评价。

子研究三：

研究前提：已经分析出系统性风险的主要成因；

研究目标：针对系统性风险的主要成因，分析在不同的成因与影响因素的状态下，在遭受相应的冲击时，监管当局应该如何决策，采取适当的政策工具或者政策工具组合，作为最优救助策略？

研究思路：介绍现有的应对系统性风险的主要的救助策略，并将其纳入待分析的系统性风险的救助策略。实验的目标是得到在不同的风险因素下、在不同的冲击影响下，哪些政策工具的单独使用和组合运用可以对不同情形下得到相对最好的效果。以上分析旨在获得最小化系统性风险的最优救助策略，并提出实现该策略的政策建议，为银行业系统性风险的防范与化解提供在所研究对象和范围内的较为全面和合理的贡献。

参数估计部分，拟考虑利用部分公开数据等作为参考，对参数进行校准和估计。估计方法拟采用Bayes方法。

\subsubsection{研究方法}

\newpage
\subsection{研究创新和不足}

\subsubsection{论文的创新}

论文创新如下：

现有文献大多对于银行业的细分程度不够。我们将对银行业从职能上（如商业银行、城商行、农信社的职能等）、规模上（四大国有银行、其余上市银行、小银行）进行细分；

在现有的DSGE框架的基础上，引入银行间市场和中央银行，引入银行间市场传染模型，构建一个复合的模型。可以综合考虑银行间的关联性在系统性风险成因的作用。

可以同时分析多种政策工具对系统性风险的效果影响，使得分析更为全面而不至于陷入局部；



\subsubsection{研究不足}

\newpage
\section{系统性风险相关理论基础}

\newpage
\subsection{相关概念的界定}

\subsubsection{系统性风险的定义}

\subsubsection{系统性风险的形成机制}

\subsubsection{系统性风险防范与化解措施}

\newpage
\subsection{系统性风险的测度}

\subsubsection{基于网络方法的测度}

\newpage
\subsection{系统性风险的形成机理}

\subsubsection{金融脆弱性理论}

\subsubsection{银行挤兑理论}

\newpage
\subsection{系统性风险的防范与化解}

\subsubsection{最后贷款人相关理论}

\subsubsection{公开市场操作相关理论}

\subsubsection{宏观审慎政策相关理论}

\newpage
\subsection{系统性风险的相关研究方法基础}

\subsubsection{复杂网络理论}

\subsubsection{银行系统基于网络的分析方法}

\subsubsection{银行系统基于网络的压力测试法}

\subsubsection{贝叶斯估计方法介绍}

\subsubsection{动态随机一般均衡模型理论介绍}

\newpage
\subsection{本章小结}

\newpage
\section{系统性风险现状分析、测度、形成机制的初步分析}

\newpage
\subsection{我国银行系统性风险的现状分析}

\newpage
\subsection{我国银行系统性风险的测度}

\subsubsection{中介指标的选取标准}

\subsubsection{中介指标的选取过程}

\subsubsection{系统性风险最终指标的构建}

\newpage
\subsection{现有系统性风险形成机制分析}

\subsubsection{信用风险的形成}

\subsubsection{信用风险导致挤兑现象对流动性风险的传导}

\subsubsection{单个金融机构流动性风险对系统性风险的传染}

\subsubsection{银行关联性对系统性风险形成的影响}

%% \subsubsubsection{3.3.4.1.	银行资产负债规模和结构}
%% \subsubsubsection{3.3.4.2.	银行间市场的网络结构}
%% \subsubsubsection{3.3.4.3.	银行间市场的网络复杂度}
%% \subsubsubsection{3.3.4.3.	银行间市场的网络复杂度}

\subsubsection{系统重要性机构对系统性风险形成的影响}

\newpage
\subsection{系统性风险的中介指标与最终指标的选取和构建}

\subsubsection{中介目标的选取标准}
\subsubsection{中介目标的选取过程}
\subsubsection{系统性风险最终指标的构建}

\newpage
\subsection{本章小结}

\newpage
\section{构筑宏观计量经济主模型框架与系统性风险形成机制的定量分析}

\newpage
\subsection{研究框架总体联系}

\subsubsection{各部门的划分}

\subsubsection{各部门的联系}
\subsubsection{各部门的行为}
\subsubsection{各部门的目标}

\newpage
\subsection{银行系统性风险影响因素选取}

\subsubsection{银行规模与银行类别}
\subsubsection{银行系统资产负债结构}
\subsubsection{银行系统资产负债结构}
\subsubsection{银行间市场}
\subsubsection{银行资产抛售规则}
\subsubsection{银行业的竞争与合作关系及程度}
\subsubsection{房地产市场相关影响因素}
\subsubsection{银行资产价格变动}

\newpage
\subsection{银行系统性风险可能的传导机制的描述}

\subsubsection{银行信用风险导致流动性风险导致银行间市场风险传染}
\subsubsection{金融资产价格波动导致抛售导致风险传染}

\newpage
\subsection{商业银行系统的异质性描述}

\subsubsection{依照商业银行的职能划分}
\subsubsection{依照商业银行的规模划分}
\subsubsection{依照商业银行之间的资产关联性划分}
\subsubsection{依照商业银行之间的同业业务关联性描述银行系统的宏观特性}
\subsubsection{依照商业银行之间的竞争与合作行为描述银行系统的宏观特性}

\newpage
\subsection{银行间市场风险传染网络模型的构建}

\subsubsection{构建银行间资产负债双边网络矩阵}
\subsubsection{构建银行间市场传染网络模型}

\newpage
\subsection{引入银行间市场网络结构到主模型框架的过程}

\newpage
\subsection{各市场的出清条件}

\newpage
\subsection{模型求解与简化}

\subsubsection{模型的稳态值求解}
\subsubsection{模型表达式的泰勒展开}

\newpage
\subsection{参数确定}

\subsubsection{数据来源与预处理}
\subsubsection{模型的识别}
\subsubsection{反映稳态特性的参数的校准}
\subsubsection{反映动态特性的参数的贝叶斯估计}

\newpage
\subsection{系统性风险成因的响应分析}

\subsubsection{外生冲击的设置}
\subsubsection{脉冲响应结果}

\newpage
\subsection{判断系统性风险成因}

\subsubsection{根据结果计算各中介指标}
\subsubsection{计算系统性风险指标}
\subsubsection{系统性风险成因判断}

\newpage
\subsection{运用结构向量自回归方法对系统性风险形成机制的分析}

\newpage
\subsection{本章小结}

\newpage
\section{银行系统性风险最优救助工具与采用的救助策略分析}

\newpage
\subsection{问题描述}

\newpage
\subsection{政策救助工具描述}

\subsubsection{注资政策}

\newpage
\subsection{模型假定}

\newpage
\subsection{模型扩建}

\subsubsection{现有的影响银行系统性风险的政策工具的选择与描述}

% 5.4.1.1.	各政策工具的特点
% 5.4.1.2.	各政策工具的优缺点
% 5.4.1.3.	各政策工具的操作行为
% 5.4.1.4.	各政策工具的传导机制

\subsubsection{银行系统性风险政策工具的引入}

% 5.4.2.1.	注资政策
% 5.4.2.2.	流动性工具
% 5.4.2.3.	银行拆分/合并
% 5.4.2.4.	公开市场操作
% 5.4.2.5.	信贷引导政策
% 5.4.2.6.	我国常用的几项宏观审慎政策

\newpage
\subsection{模型求解与简化}

\subsubsection{模型的稳态值求解}
\subsubsection{模型表达式的泰勒展开}

\newpage
\subsection{参数确定}

\subsubsection{数据来源与预处理}
\subsubsection{模型的识别}
\subsubsection{反映稳态特性的参数的校准}
\subsubsection{反映动态特性的参数的贝叶斯估计}

\newpage
\subsection{实验结果与救助策略分析}

\subsubsection{单政策工具运用时对系统性风险的影响的效果比较}
\subsubsection{政策工具组合运用时对系统性风险的影响的效果比较}

\newpage
\subsection{对上述分析结果的综合政策建议}

\newpage
\subsection{本章小结}

\newpage
\section{总结}

\newpage
\subsection{局限和不足}
\subsection{展望}

\newpage
\section{致谢}
















